% SHARD-ID: AQM-40-PRODUCT-FIELD-SPECTRUM-QUDIT-STABILIZERS
% SHARD-TITLE: Product-Field Spectra and \(n\)-Qudit Stabilizers
% SHARD-SUMMARY: Computes \(\operatorname{Spec}(k^n)\) for a finite field \(k=\mathbb F_q\).
% SHARD-SUMMARY: Proves that the residue-qudit arithmetic field is the ambient \(n\)-qudit Weyl/Pauli system.
% SHARD-SUMMARY: Separates the ambient \(n\)-qudit system from the extra choice of stabilizer labels and phases.
% SHARD-KEYWORDS: finite field, product ring, spectrum, discrete scheme, qudit, stabilizer, Weyl-Heisenberg

\section{Product-Field Spectra and n-Qudit Stabilizers}
\label{sec:product-field-spectrum-qudit-stabilizers}

\subsection{Status and Source Anchors}

This shard studies \(X=\Spec R\) for
\[
  R=k^n=k\times\cdots\times k,\qquad k=\mathbb F_q.
\]
It is a finite reduced spectrum example in the residue-qudit convention of
AQM-23, AQM-24, and AQM-28.  The Weyl operator formulas use the same finite
field sources as AQM-22 and AQM-23: Ashikhmin--Knill,
\path{references/symplectic_qecc/ashikhmin_knill_0005008_source/n9.tex},
lines 249--313 for finite-field shift and phase operators, and
Bombin--Martin--Delgado,
\path{references/toric_code/bombin_martin_delgado_0605094_source/HomologicalErrorCorrection6.tex},
lines 1208--1259 for qudit Paulis and symplectic commutators.

The new product-field convention is \texttt{CONVENTIONS.md} (ad).  The key
point is that \(\Spec(k^n)\) has \(n\) points, not \(q^n\) points.

\subsection{Lamport Dependency Skeleton}

\[
\begin{array}{rcl}
D1&:&k=\mathbb F_q,\ R=k^n,\ e_i=(0,\ldots,1,\ldots,0).\\
L1&:&\operatorname{Spec}R=\{\mathfrak m_i=\ker(\operatorname{pr}_i)\}_{i=1}^n.\\
L2&:&\kappa(\mathfrak m_i)=k,\quad D(e_i)=\{\mathfrak m_i\}.\\
D2&:&U\subset X\text{ corresponds to }I_U\subset\{1,\ldots,n\}.\\
D3&:&E(U)=\bigoplus_{i\in I_U}k^2
      =k_X^{I_U}\oplus k_Z^{I_U}.\\
D4&:&\mathcal H(U)=\bigotimes_{i\in I_U}\ell^2(k).\\
L3&:&W_U(x,z)=T_xR_z\text{ satisfies the finite-field Weyl law}.\\
T1&:&\operatorname{span}\{W_U(e):e\in E(U)\}
      =\bigotimes_{i\in I_U}M_q(\mathbb C).\\
D5&:&U\subset V\text{ gives extension by zero and tensoring by identities}.\\
T2&:&X=\Spec(k^n)\text{ gives the ambient }n\text{-qudit Weyl net}.\\
T3&:&\text{Stabilizer codes are extra isotropic labels plus phases.}
\end{array}
\]

\subsection{The Product-Field Spectrum D1--L2}

Let \(R=k^n\).  Write \(\operatorname{pr}_i:R\to k\) for the \(i\)-th
projection and \(e_i\in R\) for the idempotent with \(1\) in the \(i\)-th
coordinate and \(0\) elsewhere.

\paragraph{Lemma \(L1\).}
The prime ideals of \(R\) are exactly
\[
  \mathfrak m_i=\ker(\operatorname{pr}_i)
  =
  k\times\cdots\times k\times0\times k\times\cdots\times k .
\]
Hence \(\Spec R\) has \(n\) points.

\emph{Proof.}
First, \(R/\mathfrak m_i\cong k\), so each \(\mathfrak m_i\) is maximal and
therefore prime.

Conversely, let \(\mathfrak p\subset R\) be prime.  Since
\[
  e_ie_j=0\quad (i\ne j),\qquad 1=e_1+\cdots+e_n,
\]
at most one idempotent \(e_i\) can be outside \(\mathfrak p\): if
\(e_i,e_j\notin\mathfrak p\) for \(i\ne j\), primeness would be contradicted
by \(e_ie_j=0\in\mathfrak p\).  At least one \(e_i\) is outside
\(\mathfrak p\), because otherwise \(1=\sum_ie_i\in\mathfrak p\).  Let this
unique index be \(i\).  Then \(e_j\in\mathfrak p\) for \(j\ne i\).  Since
\(e_iR\cong k\) is a field and \(\mathfrak p\cap e_iR\) is a proper ideal, it
is \(0\).  Thus \(\mathfrak p\) is exactly the direct product of the
\(j\ne i\) factors and \(0\) in the \(i\)-th factor, namely
\(\mathfrak m_i\).

\paragraph{Lemma \(L2\).}
For every \(i\),
\[
  \kappa(\mathfrak m_i)\cong k,
  \qquad
  D(e_i)=\{\mathfrak m_i\}.
\]
Consequently \(X=\Spec R\) is a finite discrete topological space.

\emph{Proof.}
The quotient \(R/\mathfrak m_i\) is \(k\), so the residue field at
\(\mathfrak m_i\) is \(k\).  The standard open \(D(e_i)\) consists of prime
ideals not containing \(e_i\).  By the proof of Lemma \(L1\), \(e_i\notin
\mathfrak m_i\), while \(e_i\in\mathfrak m_j\) for \(j\ne i\).  Hence
\(D(e_i)=\{\mathfrak m_i\}\).  Since every singleton is open, every subset is
open.

\subsection{The Residue-Qudit Field D2--L3}

Let \(U\subset X\), and let \(I_U\subset\{1,\ldots,n\}\) be the corresponding
index set.  Define
\[
  E(U)=\bigoplus_{i\in I_U}\kappa(\mathfrak m_i)^2
      =k_X^{I_U}\oplus k_Z^{I_U}.
\]
Write \(e=(x,z)\), with \(x,z:I_U\to k\).  Let
\[
  \chi(a)=\exp\!\left(\frac{2\pi i}{p}
  \operatorname{Tr}_{k/\mathbb F_p}(a)\right)
\]
be the standard additive character.  Define
\[
  \Omega_U(e,f)=\sum_{i\in I_U}(z_ix'_i-x_iz'_i)\in k
\]
and the commutator bicharacter
\[
  B_U(e,f)=\chi(\Omega_U(e,f)).
\]

The Hilbert space is
\[
  \mathcal H(U)=\bigotimes_{i\in I_U}\ell^2(k)
  \cong \ell^2(k^{I_U}).
\]
On the basis \(|s\rangle\), \(s:I_U\to k\), define
\[
  T_x|s\rangle=|s+x\rangle,\qquad
  R_z|s\rangle=\chi\!\left(\sum_{i\in I_U}z_is_i\right)|s\rangle,
\]
and
\[
  W_U(x,z)=T_xR_z .
\]

\paragraph{Lemma \(L3\).}
For \(e=(x,z)\) and \(f=(x',z')\),
\[
  W_U(e)W_U(f)
  =
  \chi\!\left(\sum_{i\in I_U}z_ix'_i\right)W_U(e+f),
\]
and
\[
  W_U(e)W_U(f)=B_U(e,f)W_U(f)W_U(e).
\]

\emph{Proof.}
Translations add and phases add:
\[
  T_xT_{x'}=T_{x+x'},\qquad R_zR_{z'}=R_{z+z'}.
\]
Also
\[
\begin{aligned}
  R_zT_{x'}|s\rangle
  &=R_z|s+x'\rangle\\
  &=\chi\!\left(\sum_i z_i(s_i+x'_i)\right)|s+x'\rangle\\
  &=\chi\!\left(\sum_i z_ix'_i\right)T_{x'}R_z|s\rangle .
\end{aligned}
\]
Multiplying \(T_xR_zT_{x'}R_{z'}\) gives the first formula.  Dividing this
formula by the same formula with \(e\) and \(f\) exchanged gives the
commutator factor \(\chi(\Omega_U(e,f))\).

\subsection{Operator Algebra T1--T2}

\paragraph{Theorem \(T1\).}
\[
  \operatorname{span}_{\mathbb C}\{W_U(e):e\in E(U)\}
  =
  \operatorname{End}(\mathcal H(U))
  \cong
  \bigotimes_{i\in I_U}M_q(\mathbb C).
\]

\emph{Proof.}
There are \(|E(U)|=q^{2|I_U|}\) Weyl operators and
\[
  \dim_\mathbb C\operatorname{End}(\mathcal H(U))
  =
  (q^{|I_U|})^2
  =
  q^{2|I_U|}.
\]
It is enough to prove linear independence.  If \(e=(x,z)\) and
\(f=(x',z')\), the trace of \(W_U(e)^*W_U(f)\) is zero when \(x\ne x'\),
because the resulting translation has no fixed basis vector.  If \(x=x'\) but
\(z\ne z'\), the trace is
\[
  \sum_{s\in k^{I_U}}\chi\!\left(\sum_i(z'_i-z_i)s_i\right)=0,
\]
because \(z'-z\) is a nonzero linear functional and \(\chi\) is nontrivial.
If \(e=f\), the trace is \(q^{|I_U|}\).  Thus the Weyl operators are
orthogonal and nonzero for the Hilbert-Schmidt pairing, hence form a basis.
The tensor-product matrix identification follows from
\(\mathcal H(U)=\bigotimes_{i\in I_U}\ell^2(k)\).

\paragraph{Local net D5.}
If \(U\subset V\), include \(E(U)\) in \(E(V)\) by extension by zero.  The
operator inclusion is
\[
  \mathcal A(U)\longrightarrow\mathcal A(V),
  \qquad
  a\longmapsto a\otimes\mathbf 1_{I_V\setminus I_U}.
\]
The identity and composition laws are immediate from tensoring with identity
operators on the missing tensor factors.  Since \(X\) is discrete, this is the
usual tensor-factor local net over all subsets of \(\{1,\ldots,n\}\).

\paragraph{Theorem \(T2\).}
For \(X=\Spec(k^n)\), the residue-qudit arithmetic field is the ambient
standard \(n\)-qudit Weyl/Pauli net for \(q\)-level qudits:
\[
  E(X)=k^n_X\oplus k^n_Z,\qquad
  \mathcal H(X)=\ell^2(k)^{\otimes n},\qquad
  \mathcal A(X)\cong M_{q^n}(\mathbb C).
\]

\emph{Proof.}
By Lemma \(L1\), \(X\) has exactly the \(n\) points
\(\mathfrak m_1,\ldots,\mathfrak m_n\).  By Lemma \(L2\), every residue field
is \(k\), and the topology is discrete.  Therefore the full open \(X\) has
label space
\[
  E(X)=\bigoplus_{i=1}^nk^2=k^n_X\oplus k^n_Z
\]
and Hilbert space \(\bigotimes_{i=1}^n\ell^2(k)\).  Lemma \(L3\) is exactly
the finite-field multi-qudit Weyl law, and Theorem \(T1\) gives the full
matrix algebra on \(n\) \(q\)-level tensor factors.

\subsection{Stabilizer Comparison T3}

\paragraph{Theorem \(T3\).}
The ring \(k^n\) supplies the ambient \(n\)-qudit Weyl/Pauli system.  A
stabilizer code is additional data: an isotropic additive subgroup
\[
  L\leq E(X)
\]
for the bicharacter \(B_X\), together with compatible phases/eigenvalues.  If
one restricts \(L\) to be \(k\)-linear and requires \(\Omega_X(L,L)=0\) in
\(k\), one gets the usual \(k\)-linear \(q\)-ary stabilizer subclass.  Allowing
\(\mathbb F_p\)-additive \(L\) with
\[
  \operatorname{Tr}_{k/\mathbb F_p}\Omega_X(L,L)=0
\]
gives the standard finite-field additive stabilizer convention.

\emph{Proof.}
By Lemma \(L3\), two Weyl operators commute exactly when
\[
  B_X(e,f)=\chi(\Omega_X(e,f))=1.
\]
Thus an additive subgroup \(L\) gives commuting projective stabilizer labels
exactly when \(B_X(e,f)=1\) for all \(e,f\in L\), equivalently
\(\operatorname{Tr}_{k/\mathbb F_p}\Omega_X(e,f)=0\) for all \(e,f\in L\).
The phases/eigenvalues are not determined by \(R\); as in AQM-22 and AQM-38,
they are a separate lift of the commuting projective labels to actual
operators with chosen joint eigenvalues.

If \(L\) is a \(k\)-linear subspace and the stronger equality
\(\Omega_X(L,L)=0\) holds in \(k\), then the trace condition holds
automatically.  This is the \(k\)-linear stabilizer subclass.  Conversely, for
general \(q=p^r\), the trace condition may hold for an \(\mathbb F_p\)-linear
additive subgroup that is not closed under multiplication by \(k\).  Such
subgroups are part of the standard additive finite-field stabilizer formalism
but are not \(k\)-linear subspaces.  When \(q=p\) is prime, these two notions
coincide.

\subsection{Conclusion}

For \(R=k^n\), with \(X=\Spec R\), the hoped-for statement is correct with the
standard stabilizer caveat.  The product ring is the coordinate ring of a
finite discrete \(n\)-point \(k\)-scheme.  Its residue-qudit arithmetic field
is exactly the ambient \(n\) \(q\)-level qudit Weyl/Pauli system.  Stabilizer
codes do not come for free from the ring: they are later choices of isotropic
label subgroups and compatible phases.
